% concatenated from mcp-mean.tex
\documentclass[a4paper,12pt,reqno]{amsart}

\usepackage[
  a4paper,
  top=2.4cm,
  bottom=2.4cm,
  left=2.2cm,
  right=2.2cm
]{geometry}
\usepackage[british]{babel}
\usepackage[T1]{fontenc}
\usepackage{lmodern}
\usepackage[utf8]{inputenc}
\usepackage{amsmath,amssymb,amsthm,mathtools}
\usepackage[hidelinks]{hyperref}
\usepackage[capitalise,noabbrev]{cleveref}
\usepackage{setspace}
\usepackage{xcolor} 
\hypersetup{
  pdftitle={Monochromatic cycle partitions in 3-mean edge-colourings},
  pdfauthor={Richard Lang and Guilherme Oliveira Mota},
  pdfsubject={Monochromatic cycle partitions under a mean colour-degree bound},
  pdfkeywords={mean edge-colouring, local edge-colouring, monochromatic cycle,
    cycle partition}
}
 
\newtheorem{theorem}{Theorem}[section]
\newtheorem{lemma}[theorem]{Lemma}
\newtheorem{remark}[theorem]{Remark}
\newtheorem{observation}[theorem]{Observation}

\crefname{theorem}{Theorem}{Theorems}
\crefname{lemma}{Lemma}{Lemmas}
\crefname{remark}{Remark}{Remarks}
\crefname{observation}{Observation}{Observations}

\newcommand{\cP}{\mathcal P}
\newcommand{\dd}{d_{\chi}}

\title{Monochromatic cycle partitions in $3$-mean edge-colourings}
\author{Richard Lang}
\address{Departament de Matem\`atiques, Universitat Polit\`ecnica de Catalunya,
  Barcelona, Spain; and Centre de Recerca Matem\`atica, Barcelona, Spain}
\email{richard.lang@upc.edu}
\author{Guilherme Oliveira Mota}
\address{Departamento de Ci\^encia da Computa\c{c}\~ao, Instituto de Matem\'atica,
  Estat\'istica e Ci\^encia da Computa\c{c}\~ao, Universidade de S\~ao Paulo,
  Rua do Mat\~ao 1010, 05508-090 S\~ao Paulo, Brazil}
\email{mota@ime.usp.br}



\subjclass[2020]{Primary 05D10; Secondary 05C38, 05C70}
\keywords{Mean edge-colouring, local edge-colouring, monochromatic cycle,
  cycle partition}
\date{}
\onehalfspacing

\begin{document}

\begin{abstract}
  Given $r \in \mathbb{N}$ and an edge-coloured complete graph such that the average number of colours incident with a vertex is at most $r$, Conlon and Stein asked whether there is a vertex partition into a bounded
  number of monochromatic cycles, and proved this for $r=2$. We settle the
  first open case by proving the corresponding statement for $r=3$.
\end{abstract}

\maketitle

\section{Introduction}
Problems on monochromatic coverings and partitions in edge-colourings of
graphs can be traced back to a conjecture posed by Lehel in the 1970s: every
red/blue edge-colouring of the complete graph $K_n$ contains a red cycle and a
blue cycle whose vertex sets form a partition of $V(K_n)$, where edges,
vertices, and the empty set are regarded as cycles. Following earlier work of
Gyárfás~\cite{Gyarfas1983}, Łuczak, Rödl and
Szemerédi~\cite{LuczakRodlSzemeredi1998} proved this conjecture for
sufficiently large~$n$. 
Subsequently, Bessy and Thomassé~\cite{BessyThomasse2010} established the conjecture for every~$n$.

Since the 1990s, research on monochromatic cycle partitions has developed in
several directions. A prominent question asks how much of Lehel's conjecture
remains valid when more than two colours are allowed. Erdős, Gyárfás and
Pyber~\cite{ErdosGyarfasPyber1991} proved that the vertex set of every
$r$-edge-coloured complete graph can be partitioned into at most $25r^2\log r$
monochromatic cycles, and conjectured that $r$ cycles should always be enough.
For sufficiently large~$n$, Gyárfás, Ruszinkó, Sárközy and
Szemerédi~\cite{GRSS2006} improved the upper bound to $100r\log r$. The proposed
bound of $r$ cycles, however, was disproved by Pokrovskiy~\cite{Pokrovskiy2014},
who constructed, for every $r\geq 3$ and infinitely many values of $n$, an
$r$-edge-colouring of $K_n$ for which at least $r+1$ cycles are necessary.
Determining the optimal bound remains a hard problem, even for a small
number of colours. Related host graph variants replace the complete
graph by graphs satisfying minimum-degree or independence-number conditions
\cite{AllenEtAl2024,DeBiasioNelsen2017,DiBraccioPatel2026,
KorandiLangLetzterPokrovskiy2021,Letzter2019,Sarkozy2011}. Random graph
analogues were initiated by Kor\'andi, Mousset, Nenadov, \v{S}kori\'c and
Sudakov~\cite{KorandiEtAl2018} and further developed by the first author and
Lo~\cite{LangLo2021}.

Variants of covering and partitioning problems that impose conditions 
weaker than a bound on the total number of colours have also attracted
considerable attention; see the survey of Gy\'arf\'as~\cite{Gyarfas2016}. Given
an edge-colouring $\chi$ of a graph $G$, define the \emph{palette} of $v$ as
  \[
  \cP(v):=\{\chi(vu)\colon vu\in E(G)\}\,,
  \]
and let $\dd(v):=|\cP(v)|$ be the \emph{colour-degree} of $v$. We could allow an
unbounded number of colours, while restricting the number of colours on edges
incident to each vertex as in the following definition: a colouring $\chi$ is
\emph{$r$-local} if $\dd(v)\le r$ for every vertex $v$. Even though the number
of colours can be unbounded, the local hypothesis of an $r$-local colouring
provides a uniform bound for the colour-degree at every vertex.

Conlon and Stein \cite{ConlonStein2016} proved that every $r$-locally coloured
complete graph has a partition into $O(r^2\log r)$ monochromatic cycles; for
$r=2$, they showed that two cycles of different colours suffice. For
sufficiently large graphs, the first author and Stein~\cite{LangStein2017}
obtained the explicit bounds $2r^2$ for $r$-locally coloured complete graphs
and $4r^2$ for $r$-locally coloured balanced complete bipartite graphs.
Complete bipartite host graphs have also been studied under global colour
bounds~\cite{BenevidesQuintinoTalon2025,Haxell1997,
LangSchaudtStein2017}. For sufficiently large complete graphs,
S\'ark\"ozy~\cite{Sarkozy2020} subsequently improved the dependence for
locally coloured complete graphs to $O(r\log r)$.

Relaxing the condition on the colouring even more, one could allow a small
exceptional set of vertices to have large palettes, which is not the
case in $r$-local colourings. This motivates the definition of an
\emph{$r$-mean} colouring of an $n$-vertex graph $G$, which is a colouring satisfying 
\[
  \frac{\sum_{v\in V(G)}\dd(v)}{n} \le r\,.
\]
Mean colourings arose in the Ramsey-theoretic work of Caro \cite{Caro1992} and
were subsequently studied systematically by Caro and Tuza \cite{CaroTuza1993}.
Conlon and Stein also initiated the monochromatic cycle-partition problem for
mean colourings in complete graphs. They proved that every $2$-mean coloured
complete graph can be partitioned into two monochromatic cycles of different
colours, and asked whether, for every fixed $r$, there is a constant $C$ such
that every $r$-mean coloured complete graph has a partition into at most $C$
monochromatic cycles. We settle the first open case by proving the analogue of
Conlon and Stein's result for $r=3$. 

\begin{theorem}\label{thm:main}
  There is an absolute constant $C$ such that the vertex set of every
  $3$-mean edge-coloured complete graph can be partitioned into at most $C$
  disjoint monochromatic cycles.
\end{theorem}

Our argument relies on structural properties specific to the case $r=3$, and
genuinely new ideas seem to be required to deal with larger values of $r$.
We finish this introduction by presenting an overview of the proof of
\Cref{thm:main}. In \Cref{sec:preliminaries} we collect the three covering
results used in the argument and prove a one-sided bipartite lemma that absorbs
the part of the vertex set that is dense in a common colour, and
in~\Cref{sec:proof} we prove~\Cref{thm:main}.

\vspace{0.2cm}
\noindent\textit{Overview of the proof.} To prove \cref{thm:main}, we partition the vertices into a low
colour-degree set $A$, an intermediate set $M$, and a high colour-degree set
$B$, and write $a:=|A|$ and $b:=|B|$.  The set $M$ is immediately covered by a
bounded number of monochromatic
cycles from a result by Conlon and Stein (\Cref{thm:local-complete}). The mean
inequality gives~${a\ge6b}$. Moreover, the palettes in $A$ form a
pairwise-intersecting family of sets of size at most two. Consequently, either
the union of these palettes has size at most three, or they all contain a common colour~$c$
(\Cref{lem:palette-dichotomy}). In the first case, if $b$ is not sufficiently
large for \Cref{thm:three-colour-bipartite}, we cover $B$ by singleton cycles;
otherwise, any $b$ vertices of $A$ form with $B$ a balanced bipartite graph
using at most three colours.

Consider the second case and let $A_c$ be the set of vertices of $A$ whose
palette contains only colour $c$, and put $a_c:=|A_c|$. If $a_c\ge b$, then $B$
can be paired with $b$ vertices of $A_c$ in colour $c$. Otherwise, the mean
inequality allows us first to pair all of $A_c$ with an equally large set
$B_0\subseteq B$ in colour $c$. Put $X:=A\setminus A_c$ and $Y:=B\setminus B_0$.
Then $|X|\ge12|Y|$, and every vertex of $X$ has colour-degree exactly two. We
split the vertices of $Y$ according to their colour-$c$ degrees into $X$. The
vertices having at least $3|X|/8$ colour-$c$ neighbours are absorbed by at most
three colour-$c$ cycles, using no more vertices of $X$ than vertices already
covered (\Cref{lem:dense-common}, proved using \Cref{lem:posa}). Let $Y'$ be the
set left after this step and let~$X_0\subseteq X$ be the \emph{reservoir}
vertices used by the absorbing cycles.

Every vertex of $Y'$ has more than $5|X|/8$ non-$c$ neighbours in $X$. If $|Y'|$
is not sufficiently large for \Cref{thm:three-colour-bipartite}, we cover it by
singleton cycles. Otherwise, the residual colour-degree budget allows us to
apply \Cref{lem:bounded-colour-core}. Its proof uses one witness vertex of $Y'$ to
choose the partner set inside at most two non-$c$ colour classes. A non-$c$
witness edge at a vertex~$u\in X$ forces its palette to consist of $c$ and the
colour of that edge. Hence the resulting balanced bipartite graph uses at most
three colours, and a result obtained by the first author, Schaudt and Stein
(\Cref{thm:three-colour-bipartite}) partitions its vertices into a bounded
number of monochromatic cycles. Finally, we are able to cover all uncovered
vertices by using \Cref{thm:local-complete} as these vertices have colour-degree
at most two.


\section{Preliminaries}\label{sec:preliminaries}

We use three known results at different stages of the proof.  P\'osa's lemma
converts a bound on the independence number of an auxiliary graph into a cycle
partition; we shall use it to prove our one-sided bipartite covering lemma.  The
complete graph theorem of Conlon and Stein serves as a cleaning tool: it covers
both the vertices of intermediate colour-degree and the final residue of
colour-degree at most two.  The balanced bipartite theorem of the first author,
Schaudt and Stein is used whenever we obtain a sufficiently large balanced
bipartite graph using at most three colours: directly in the first case of the
palette dichotomy, and again after constructing a partner set for the remaining
high colour-degree vertices.

\begin{lemma}[P\'osa \cite{Posa1963}]\label{lem:posa}
  The vertex set of any graph $H$ can be partitioned into at most
  $\alpha(H)$ disjoint cycles.
\end{lemma}

The next result is our complete graph cleaning tool. Notice that passing to
an induced subgraph cannot increase any colour-degree, so a class defined by
an upper bound on its colour-degrees automatically satisfies the required local
hypothesis.

\begin{theorem}[Conlon--Stein \cite{ConlonStein2016}]
  \label{thm:local-complete}
  For every positive integer $r$, there is a constant $g(r)$ such that the
  vertex set of every $r$-locally coloured complete graph can be partitioned
  into at most~$g(r)$ disjoint monochromatic cycles.
\end{theorem}

We remark that one may take $g(r)=O(r^2\log r)$ in \cref{thm:local-complete}. In
the two branches of the main proof involving high colour-degree vertices, we
shall form a balanced bipartite graph whose cross-edges use at most three
colours.  The following result will cover the resulting pair.

\begin{theorem}[Lang--Schaudt--Stein \cite{LangSchaudtStein2017}]
  \label{thm:three-colour-bipartite}
  For sufficiently large $n$, the vertices of every $K_{n,n}$ whose edges are
  coloured with at most three colours can be partitioned into at most $18$
  disjoint monochromatic cycles.
\end{theorem}

The next lemma is the interface between the dense monochromatic part of the
argument and P\'osa's lemma.  It is deliberately one-sided: only the smaller
part must be covered.
Its proof is adapted from the work of Erdős, Gyárfás and Pyber~\cite{ErdosGyarfasPyber1991}.
In the main proof, $H$ will consist of the edges of one
fixed colour.

\begin{lemma}\label{lem:dense-common}
  Let $H$ be a bipartite graph with parts $X$ and $Y$, and put
  $x:=|X|$ and $y:=|Y|$. If $x\ge12y$ and $d_H(v)\ge3x/8$ for every
  $v\in Y$, then $Y$ can be covered by at most three vertex-disjoint cycles.
\end{lemma}

\begin{proof}
  The assertion is immediate when $y\le3$, so suppose that $y\ge4$ and let
  $N_H(u,v)$ denote the joint neighbourhood of vertices $u$ and $v$ in $H$. In
  order to apply~\Cref{lem:posa}, define an auxiliary graph $J$ on $Y$ by
  putting an edge $uv$ in $J$ if and only if $|N_H(u,v)|\ge y$.
  
  We claim that $\alpha(J)\le3$. Indeed, take four vertices
  $v_1,v_2,v_3,v_4\in Y$. For $w\in X$, let $t_w$ be the number of these four
  vertices adjacent to $w$. Since $\binom{t}{2}\ge t-1$ for every
  non-negative integer $t$, we have
  \[
    \sum_{1\le i<j\le4}|N_H(v_i,v_j)|
      =\sum_{w\in X}\binom{t_w}{2}
      \ge\sum_{w\in X}(t_w-1)
      \ge\frac{3x}{2}-x
      =\frac{x}{2}\ge6y\,.
  \]
  Thus some pair has at least $y$ common neighbours and is an edge of $J$.
  Every four vertices of $J$ therefore span an edge, proving the claim.

  By \cref{lem:posa}, the vertices of $J$ can be partitioned into at most three
  cycles. We now realise these auxiliary cycles as cycles in $H$. For every pair
  $u,v$ of consecutive vertices on an auxiliary cycle of order at least three,
  request one vertex $z_{uv}$ from $N_H(u,v)$. If an auxiliary cycle has order
  two, and hence consists of a single edge $uv$, request two distinct vertices
  $z_1,z_2$ from $N_H(u,v)$. No vertex is requested for a singleton cycle. The
  total number of requests is the sum of the orders of the non-singleton
  auxiliary cycles, so it is at most $y$. Each requested common neighbourhood
  has size at least $y$ by the definition of $J$ and we may therefore choose
  all the requested vertices to be distinct.

  For an auxiliary cycle of order at least three in $J$, replace each of its
  edges $uv$ by the path~$u z_{uv}v$, giving a cycle in $H$. For an
  auxiliary cycle of order two, the four-vertex cycle~$uz_1vz_2u$ is a cycle of
  $H$. Finally, retain every singleton cycle as a singleton in $H$. Since all
  the chosen vertices are distinct, this yields at most three vertex-disjoint
  cycles of $H$ covering~$Y$, and it uses at most $y$ vertices of $X$.
\end{proof}

\section{Proof of the main theorem}\label{sec:proof}

We first isolate the two ingredients that drive the proof. Given an
edge-coloured graph $G$ and a vertex set $S\subseteq V(G)$, define
$\dd(S):=\sum_{v\in S}\dd(v)$.


\begin{lemma}[Palette dichotomy]\label{lem:palette-dichotomy} Let $H$ be an
edge-coloured complete graph, and let $A$ and~$B\neq\emptyset$ be disjoint
subsets of $V(H)$. If $\dd(A)+\dd(B)\le3(|A|+|B|)$ and
\begin{align*}
&\dd(u) \leq 2\text{ for every }u\in A,\quad\text{and}\\
&\dd(v) \geq 15 \text{ for every }v\in B\,,
\end{align*}
then $|A| \ge 6|B|$ and either all palettes of vertices in $A$ contain a common
colour, or
\[
  \left|\bigcup_{u\in A}\cP(u)\right|\le3.
\]
\end{lemma}

\begin{proof}
Put $a:=|A|$ and $b:=|B|$. Since every palette is non-empty,
  \[
    a+15b\le \dd(A)+\dd(B)\le3(a+b)\,,
  \]
  and hence $a\ge6b\ge6$. For any two $u,v\in A$, the colour of $uv$
  belongs to $\cP(u)\cap\cP(v)$. Thus the palettes in $A$ form a
  pairwise-intersecting family of non-empty sets of size at most two.
  Suppose that this family has no common colour. It contains no singleton, since
  the colour in a singleton would then belong to every palette. Choose a palette
  $\{\alpha,\beta\}$. Since neither~$\alpha$ nor~$\beta$ is common, pairwise
  intersection gives palettes $\{\beta,\gamma\}$ and $\{\alpha,\gamma\}$ for
  some third colour~$\gamma$. Every set of size two meeting all three of these
  palettes is one of $\{\alpha,\beta\}$, $\{\alpha,\gamma\}$ and
  $\{\beta,\gamma\}$. Hence every palette in $A$ is contained in
  $\{\alpha,\beta,\gamma\}$, proving the second alternative.
\end{proof}

For an edge-coloured graph $G$, a colour $\gamma$, a vertex $v\in V(G)$, and a
set $S\subseteq V(G)\setminus\{v\}$, we write $N_\gamma(v,S)$ for the set of
vertices $u\in S$ for which $uv$ has colour $\gamma$. The next lemma
converts a bound on the total palette size of one side into the balanced
bounded-colour core required by \cref{thm:three-colour-bipartite}.

\begin{lemma}[Bounded-colour core]\label{lem:bounded-colour-core} Let $H$ be an
  edge-coloured complete graph, and let~$X$ and~$Y'\neq\emptyset$ be disjoint
subsets of $V(H)$. Let $c$ be a colour, $X_0\subseteq X$ and put~$\ell:=\lfloor 5|X|/8\rfloor+1$. If
\begin{align*}
&c\in\cP(u)\text{ and }\dd(u)\le2
  &&\text{for every }u\in X,\quad\text{and}\\
&|X\setminus N_c(v,X)|\ge\ell
  &&\text{for every }v\in Y'\,,
\end{align*}
and
\[
  15|Y'|\le\dd(Y')\le |X|+3|Y'|-12|X_0|\,,
\]
then there is a set $X^*\subseteq X\setminus X_0$ with $|X^*|=|Y'|$ such that
the complete bipartite graph between $X^*$ and $Y'$ uses at most three
colours.
\end{lemma} 

\begin{proof}
Put $x:=|X|$, $y':=|Y'|$, $x_0:=|X_0|$ and $t:=\dd(Y')/y'\geq 15$. For each
$v\in Y'$, we define $t_v:=|\cP(v)\setminus\{c\}|$. Thus, since every vertex of
$Y'$ has a non-$c$ neighbour in $X$, we have $t_v\ge1$, which combined with
$\sum_{v\in Y'}t_v\le\dd(Y')=ty'$ implies that some vertex $v\in Y'$ satisfies
$1\le t_v\le t$.

Let $\sigma:=\ell-x_0$. Deleting $X_0$ removes at most $x_0$ non-$c$ neighbours
of $v$, so $v$ has at least~$\sigma$ non-$c$ neighbours in $X\setminus X_0$.
The key estimate that will allow us to choose the $y'$ vertices forming $X^*$
is
\begin{equation}\label{eq:core-target}
  \frac{2\sigma}{t}>y'\,.
\end{equation}
To see that \eqref{eq:core-target} holds, note first that the bounds on
$\dd(Y')=ty'$ give
$$
x\ge ty'-3y'+12x_0\ge 4ty'/5+12x_0,
$$
where we used that $t\geq 15$ in the last inequality. Thus, since
$\ell>5x/8$, it follows that
$$2\sigma>5x/4-2x_0\ge ty'+13x_0\ge ty',
$$
which proves \eqref{eq:core-target}.

Now let us conclude the proof of the lemma from~\eqref{eq:core-target}. Suppose
first that $t_v=1$. Then the at least $\sigma$ available non-$c$ neighbours
of $v$ belong to the same non-$c$ colour class. By \eqref{eq:core-target} and~$t\ge15$, we have $\sigma>ty'/2>y'$, so this class contains at least $y'$
available vertices. In the case where $t_v\ge2$, write
$d_1\ge\cdots\ge d_{t_v}\ge0$ for
the orders of the corresponding non-$c$ colour classes among these neighbours,
including empty classes if necessary. Since their total order is at least
$\sigma$, we have
\[
  d_1+d_2
  \ge\frac{2}{t_v}\sum_{i=1}^{t_v}d_i
  \ge\frac{2\sigma}{t_v}
  \ge\frac{2\sigma}{t}
  >y'\,,
\]
where the penultimate inequality uses $t_v\le t$. Thus, in either case, the
union of at most two such classes contains at least $y'$ available vertices.

Choose a set $X^*\subseteq X\setminus X_0$ of order $y'$ from these classes.
Let $c_1$ and $c_2$ be their colours; if there is only one such colour, set
$c_2:=c_1$. Thus, for every $u\in X^*$, the edge $uv$ has colour $c_1$ or
$c_2$. This edge is not of colour $c$. Since $c\in\cP(u)$ and
$\dd(u)\le2$, we have
$\cP(u)\subseteq\{c,c_1,c_2\}$. Therefore every edge between $X^*$ and $Y'$
has colour in the set $\{c,c_1,c_2\}$, which has size at most three.
\end{proof}


\subsection{Proof of the main result}

We now combine the ingredients we obtained so far to prove~\Cref{thm:main}. After covering the vertices of intermediate colour-degree, we use the mean condition and the palette dichotomy to analyse the interaction between the low and high colour-degree classes. In the common-colour case, Lemmas \ref{lem:dense-common} and \ref{lem:bounded-colour-core} cover the high colour-degree
vertices, while the remaining $2$-locally coloured vertices are handled by
\Cref{thm:local-complete}.

\begin{proof}[Proof of \Cref{thm:main}]
Let $K_n$ be edge-coloured by $\chi$ and suppose that $\sum_{v\in V(K_n)}\dd(v)\le3n$. The case $n=1$ is trivial, so assume $n\ge2$.

For clarity, we divide the proof into four parts: \emph{Cleaning and
dichotomy}, \emph{Common-colour absorption}, \emph{Residual core}, and
\emph{Completion}. The first part provides
the basic reduction, the next two cover the high colour-degree vertices, and the
last one covers the remaining vertices and counts the cycles.

\medskip
\noindent\emph{Cleaning and dichotomy.} We split the vertices of $K_n$ into a
$2$-local reservoir, an immediately coverable middle class, and a high
colour-degree exceptional set as follows:
\begin{align*}
  A&:=\{v\colon\dd(v)\le2\}\,,\\
  M&:=\{v\colon3\le\dd(v)\le14\}\,,\\
  B&:=\{v\colon\dd(v)\ge15\}\,.
\end{align*}
If $M$ is non-empty, the complete graph induced by $M$ is $14$-locally
coloured, so \cref{thm:local-complete} partitions $M$ into at most $g(14)$
monochromatic cycles; if $M$ is empty, no cycle is needed.  It remains to
cover $A\cup B$.

Put $a:=|A|$ and $b:=|B|$.  Since every vertex of $M$ has colour-degree at
least three and the colouring is $3$-mean,
\begin{equation}\label{eq:residual-mean}
  \dd(A)+\dd(B)
  \le 3(a+b)\,.
\end{equation}
If $B=\varnothing$, apply \cref{thm:local-complete} to the $2$-locally coloured
complete graph induced by $A$ (possibly empty, when $A=\varnothing$).  We may
thus assume that $b>0$.

Since we have \eqref{eq:residual-mean}, we may
apply~\cref{lem:palette-dichotomy} to $A$ and $B$, obtaining $a\ge6b$.
Moreover, either 
\begin{enumerate}
\item $|\bigcup_{u\in A}\cP(u)|\le3$, or
\item all the palettes of vertices in $A$ have a common colour.
\end{enumerate}

\vspace{0.2cm}
\noindent\textit{Case 1}: $|\bigcup_{u\in A}\cP(u)|\le3$.

If $b$ is not sufficiently large for \cref{thm:three-colour-bipartite}, cover
the vertices of $B$ by singleton cycles; this uses only a bounded number of
cycles. Otherwise, choose any set $A^*\subseteq A$ of order $b$, which is
possible because $a\ge6b$. The edges between $A^*$ and $B$ use at most three
colours, so \cref{thm:three-colour-bipartite} partitions $A^*\cup B$ into at
most $18$ monochromatic cycles. In either case, all uncovered vertices lie in
$A$ and are partitioned into at most $g(2)$ monochromatic cycles by
\cref{thm:local-complete}. This proves the theorem in this case.


\vspace{0.2cm}
\noindent\textit{Case 2}: all the palettes of vertices in $A$ have a common colour.

Fix such a common colour $c$, and put
\[
  A_c:=\{u\in A\colon\cP(u)=\{c\}\},\qquad a_c:=|A_c|\,.
\]
We start with the following simple observation.
\begin{observation}\label{obs:monochromatic-biclique}
If $U\subseteq A$ and $V\subseteq B$ are non-empty sets with $|U|=|V|$ and
every edge between $U$ and $V$ has colour $c$, then $U\cup V$ is covered by one
colour-$c$ cycle (or edge).
\end{observation}

If $a_c\ge b$, we choose a set $A_c^*\subseteq A_c$ of order $b$. By
\cref{obs:monochromatic-biclique}, $A_c^*\cup B$ is covered by one colour-$c$
cycle. The remaining vertices lie in $A$ and are handled by
\cref{thm:local-complete}. We may then assume that
\begin{equation}\label{eq:a-c-small}
  a_c<b.
\end{equation}
Since every palette in $A$ contains $c$, we have
\[
  \dd(A)=a_c+2(a-a_c)=2a-a_c\,.
\]
Combining this identity, \eqref{eq:residual-mean}, and
$\dd(B)\ge15b$, we obtain
\[
  a+a_c\ge12b.
\]
Now choose a set $B_0\subseteq B$ of order $a_c$. The set $A_c\cup B_0$ can be
covered with at most one cycle. Indeed, if $a_c>0$,
\cref{obs:monochromatic-biclique} covers $A_c\cup B_0$ with one colour-$c$
cycle; if $a_c=0$, then $A_c=B_0=\varnothing$, so no cycle is needed to cover
these sets. Therefore, it remains to cover 
\[
  X:=A\setminus A_c,\qquad Y:=B\setminus B_0\,,
\]
and we write $x:=|X|$ and $y:=|Y|$. By~\eqref{eq:a-c-small}, we have $y>0$.
Moreover,
\begin{equation}\label{eq:peeled-ratio}
  x-12y=(a+a_c-12b)+10a_c\ge0\,.
\end{equation}

In summary, from now on we may assume that all the palettes of vertices in $A$
have a common colour and we have obtained disjoint sets $X,Y$ with $y>0$,
$x\ge12y$, and $\dd(u)=2$ and $c\in\cP(u)$ for every $u\in X$. The aim now is to
cover $Y$ using some vertices of $X$, with a bounded number of monochromatic
cycles, since the unused vertices of $X$ can then be handled by
\cref{thm:local-complete}.

\medskip
\noindent\emph{Common-colour absorption.} The purpose of this step is to cover
the vertices of $Y$ with many colour-$c$ neighbours in $X$, while using
only a controlled number of vertices of $X$. Every vertex left in $Y$ will
consequently have many non-$c$ neighbours in $X$, which is precisely the
condition needed to construct the bounded-colour core in the next part.

Split $Y$ into the sets
\begin{align*}
  Y':=\left\{v\in Y\colon|N_c(v,X)|<\frac{3x}{8}\right\},\quad\text{ and }\quad Y_c:=Y\setminus Y'
\end{align*}
and write $y':=|Y'|$ and $y_c:=|Y_c|=y-y'$.

The next step is to show that we can cover $Y_c$ with monochromatic cycles by
using at most~$y_c$ vertices of $X$. If $Y_c=\varnothing$, set
$X_0:=\varnothing$ and $x_0:=0$. If $Y_c$ is non-empty, we want to apply
\cref{lem:dense-common} to the bipartite graph of the colour-$c$ edges between
$X$ and $Y_c$. To see that we can apply this lemma, note that
\eqref{eq:peeled-ratio} gives $x\ge12y\ge12y_c$, and every vertex of $Y_c$ has
colour-$c$ degree at least $3x/8$ into $X$, so the conditions to apply
\cref{lem:dense-common} are satisfied. Therefore, we conclude that there are at
most three disjoint colour-$c$ cycles covering~$Y_c$. Let~$X_0\subseteq X$ be
the set of vertices of $X$ used by these cycles, and put $x_0:=|X_0|$. The last
assertion of \cref{lem:dense-common} gives $x_0\le y_c$.

Thus, in either case, $Y_c$ is covered and the set $X_0$ of vertices used from
$X$ satisfies $x_0\le y_c$; consequently, $X\setminus X_0$ remains available
for the residual core.

If $y'=0$, all remaining vertices lie in $X$ and can be handled by
\cref{thm:local-complete}. If $y'>0$ but is not sufficiently large for
\cref{thm:three-colour-bipartite}, we cover the vertices of $Y'$ by singleton
cycles and again apply \cref{thm:local-complete} to the remaining part of $X$.
We may therefore assume that $y'$ is sufficiently large for
\cref{thm:three-colour-bipartite}.

At this point, the only uncovered vertices outside $X$ are the vertices of $Y'$, and
$y'$ is sufficiently large. It remains to find a set
$X^*\subseteq X\setminus X_0$ of order $y'$ such that the bipartite graph
between $X^*$ and $Y'$ uses at most three colours; then
\cref{thm:three-colour-bipartite} covers $X^*\cup Y'$, and the remaining
vertices of $X$ can be handled by \cref{thm:local-complete}.

\medskip
\noindent\emph{Residual core.} The purpose of this step is to construct the set
$X^*$ announced above. We first show that every vertex of $Y'$ has many non-$c$
neighbours in $X$ and use the mean colour-degree condition to obtain the upper
bound on $\dd(Y')$ required by \cref{lem:bounded-colour-core}, which yields a
balanced bipartite core $X^*\cup Y'$ using at most three colours, to which
\cref{thm:three-colour-bipartite} applies.

Set $\ell:=\left\lfloor\frac{5x}{8}\right\rfloor+1$. Since
$Y':=\{v\in Y\colon|N_c(v,X)|<3x / 8\}$, for every $v\in Y'$ we have
\[
  x-|N_c(v,X)|
  \ge \left\lfloor\frac{5x}{8}\right\rfloor+1=\ell\,.
\]
This verifies the neighbourhood hypothesis of \cref{lem:bounded-colour-core}.
Indeed, after the previously used set $X_0$ is removed, every vertex of $Y'$
still has at least $\ell-x_0$ non-$c$ neighbours in $X\setminus X_0$; these are
the available vertices from which \cref{lem:bounded-colour-core} constructs
$X^*$.

Now, since every vertex of $Y'$ has colour-degree at least $15$,
\begin{equation}\label{eq:yprime-degree-lower}
  \dd(Y')\ge15y'\,.
\end{equation}
On the other hand, using \eqref{eq:residual-mean},
the identity $\dd(A)=2a-a_c$ and the inequalities $\dd(B_0)\ge15a_c$ and
$\dd(Y_c)\ge15y_c$, we have
\begin{equation}\label{eq:yprime-degree-upper}
\begin{aligned}
  \dd(Y')
  &\le3(a+b)-(2a-a_c)-15a_c-15y_c\\
  &=x+3y'-12y_c-10a_c\\
  &\le x+3y'-12x_0\,,
\end{aligned}
\end{equation}
where the last inequality follows from $x_0\le y_c$ and $a_c\ge0$.

Let us explicitly check that \cref{lem:bounded-colour-core} applies. The sets
$X$ and $Y'$ are disjoint, $Y'$ is non-empty by assumption, and
$X_0\subseteq X$. Every $u\in X$ satisfies
$c\in\cP(u)$ and $\dd(u)=2$, while the neighbourhood estimate above gives
$|X\setminus N_c(v,X)|\ge\ell$ for every $v\in Y'$. Finally, from
\eqref{eq:yprime-degree-lower} and \eqref{eq:yprime-degree-upper} we have
\[
  15y'\le\dd(Y')\le x+3y'-12x_0\,.
\]
Thus every hypothesis of \cref{lem:bounded-colour-core} has already been
established. Applying it, we obtain a set $X^*\subseteq X\setminus X_0$ of order
$y'$ such that the complete bipartite graph between $X^*$ and $Y'$ uses at most
three colours. Since $y'$ is sufficiently large,
\cref{thm:three-colour-bipartite} covers its vertices with at most $18$
disjoint monochromatic cycles.

\medskip
\noindent\emph{Completion.}
At this point every uncovered vertex lies in $X$.  The complete graph induced
by those vertices, if non-empty, is still $2$-locally coloured, so
\cref{thm:local-complete} covers it with at most~$g(2)$ disjoint
monochromatic cycles.  All cycle families constructed above are
vertex-disjoint: they successively use vertices from
\[
  M,\qquad A_c\cup B_0,\qquad Y_c\cup X_0,\qquad Y'\cup X^*,
  \qquad\text{and the remaining part of }X.
\]
The small exceptional sets treated with singleton cycles have bounded order,
and all other steps use at most $g(14)+4+g(2)+18$ cycles. Hence, in every case,
the total number of monochromatic cycles used to partition $V(K_n)$ is bounded
by an absolute constant. 
This proves \Cref{thm:main}.
\end{proof}


	\section*{Acknowledgments}

R. Lang was supported by the Ramón y Cajal programme (RYC2022-038372-I) and by
grant PID2023-147202NB-I00 funded by MICIU/AEI/10.13039/501100011033. G. O. Mota
was supported by CNPq (420838/2025-2 and 315916/2023-0) and FAPESP (2024/13859-4
and 2023/03167-5).


\begin{thebibliography}{99}

\bibitem{AllenEtAl2024}
P.~Allen, J.~B\"ottcher, R.~Lang, J.~Skokan and M.~Stein,
\newblock Partitioning a $2$-edge-coloured graph of minimum degree
  $2n/3+o(n)$ into three monochromatic cycles,
\newblock \emph{European J. Combin.} \textbf{121} (2024), Article 103838,
\newblock \href{https://doi.org/10.1016/j.ejc.2023.103838}{doi:10.1016/j.ejc.2023.103838}.

\bibitem{BenevidesQuintinoTalon2025}
F.~S.~Benevides, A.~L.~Quintino and A.~Talon,
\newblock Partitioning $2$-edge-coloured bipartite graphs into monochromatic
  cycles,
\newblock \emph{European J. Combin.} \textbf{130} (2025), Article 104192,
\newblock \href{https://doi.org/10.1016/j.ejc.2025.104192}{doi:10.1016/j.ejc.2025.104192}.

\bibitem{BessyThomasse2010}
S.~Bessy and S.~Thomass\'e,
\newblock Partitioning a graph into a cycle and an anticycle, a proof of
  Lehel's conjecture,
\newblock \emph{J. Combin. Theory Ser. B} \textbf{100} (2010), 176--180,
\newblock \href{https://doi.org/10.1016/j.jctb.2009.07.001}{doi:10.1016/j.jctb.2009.07.001}.

\bibitem{Caro1992}
Y.~Caro,
\newblock On several variations of the Tur\'{a}n and Ramsey numbers,
\newblock \emph{J. Graph Theory} \textbf{16} (1992), 257--266,
\newblock \href{https://doi.org/10.1002/jgt.3190160309}{doi:10.1002/jgt.3190160309}.

\bibitem{CaroTuza1993}
Y.~Caro and Zs.~Tuza,
\newblock On $k$-local and $k$-mean colorings of graphs and hypergraphs,
\newblock \emph{Quart. J. Math. Oxford Ser. (2)} \textbf{44} (1993), 385--398,
\newblock \href{https://doi.org/10.1093/qmath/44.4.385}{doi:10.1093/qmath/44.4.385}.

\bibitem{ConlonStein2016}
D.~Conlon and M.~Stein,
\newblock Monochromatic cycle partitions in local edge colorings,
\newblock \emph{J. Graph Theory} \textbf{81} (2016), 134--145,
\newblock \href{https://doi.org/10.1002/jgt.21867}{doi:10.1002/jgt.21867}.

\bibitem{DeBiasioNelsen2017}
L.~DeBiasio and L.~L.~Nelsen,
\newblock Monochromatic cycle partitions of graphs with large minimum degree,
\newblock \emph{J. Combin. Theory Ser. B} \textbf{122} (2017), 634--667,
\newblock \href{https://doi.org/10.1016/j.jctb.2016.08.006}{doi:10.1016/j.jctb.2016.08.006}.

\bibitem{DiBraccioPatel2026}
F.~Di Braccio and V.~Patel,
\newblock Monochromatic cycle partitions of $r$-edge-coloured graphs with
  high minimum degree,
\newblock arXiv:2601.22117 (2026),
\newblock \href{https://arxiv.org/abs/2601.22117}{arXiv:2601.22117}.

\bibitem{ErdosGyarfasPyber1991}
P.~Erd\H{o}s, A.~Gy\'arf\'as and L.~Pyber,
\newblock Vertex coverings by monochromatic cycles and trees,
\newblock \emph{J. Combin. Theory Ser. B} \textbf{51} (1991), 90--95,
\newblock \href{https://doi.org/10.1016/0095-8956(91)90007-7}{doi:10.1016/0095-8956(91)90007-7}.

\bibitem{GRSS2006}
A.~Gy\'arf\'as, M.~Ruszink\'o, G.~N.~S\'ark\"ozy and E.~Szemer\'edi,
\newblock An improved bound for the monochromatic cycle partition number,
\newblock \emph{J. Combin. Theory Ser. B} \textbf{96} (2006), 855--873,
\newblock \href{https://doi.org/10.1016/j.jctb.2006.02.007}{doi:10.1016/j.jctb.2006.02.007}.


\bibitem{Gyarfas1983}
A.~Gy\'arf\'as,
\newblock Vertex coverings by monochromatic paths and cycles,
\newblock \emph{J. Graph Theory} \textbf{7} (1983), 131--135,
\newblock \href{https://doi.org/10.1002/jgt.3190070116}{doi:10.1002/jgt.3190070116}.

\bibitem{Gyarfas2016}
A.~Gy\'arf\'as,
\newblock Vertex covers by monochromatic pieces---a survey of results and
  problems,
\newblock \emph{Discrete Math.} \textbf{339} (2016), 1970--1977,
\newblock \href{https://doi.org/10.1016/j.disc.2015.07.007}{doi:10.1016/j.disc.2015.07.007}.

\bibitem{Haxell1997}
P.~E.~Haxell,
\newblock Partitioning complete bipartite graphs by monochromatic cycles,
\newblock \emph{J. Combin. Theory Ser. B} \textbf{69} (1997), 210--218,
\newblock \href{https://doi.org/10.1006/jctb.1997.1737}{doi:10.1006/jctb.1997.1737}.

\bibitem{KorandiEtAl2018}
D.~Kor\'andi, F.~Mousset, R.~Nenadov, N.~\v{S}kori\'c and B.~Sudakov,
\newblock Monochromatic cycle covers in random graphs,
\newblock \emph{Random Structures Algorithms} \textbf{53} (2018), 667--691,
\newblock \href{https://doi.org/10.1002/rsa.20819}{doi:10.1002/rsa.20819}.

\bibitem{KorandiLangLetzterPokrovskiy2021}
D.~Kor\'andi, R.~Lang, S.~Letzter and A.~Pokrovskiy,
\newblock Minimum degree conditions for monochromatic cycle partitioning,
\newblock \emph{J. Combin. Theory Ser. B} \textbf{146} (2021), 96--123,
\newblock \href{https://doi.org/10.1016/j.jctb.2020.07.005}{doi:10.1016/j.jctb.2020.07.005}.

\bibitem{LangLo2021}
R.~Lang and A.~Lo,
\newblock Monochromatic cycle partitions in random graphs,
\newblock \emph{Combin. Probab. Comput.} \textbf{30} (2021), 136--152,
\newblock \href{https://doi.org/10.1017/S0963548320000401}{doi:10.1017/S0963548320000401}.

\bibitem{LangSchaudtStein2017}
R.~Lang, O.~Schaudt and M.~Stein,
\newblock Almost partitioning a $3$-edge-colored $K_{n,n}$ into five
  monochromatic cycles,
\newblock \emph{SIAM J. Discrete Math.} \textbf{31} (2017), 1374--1402,
\newblock \href{https://doi.org/10.1137/15M104222X}{doi:10.1137/15M104222X}.

\bibitem{LangStein2017}
R.~Lang and M.~Stein,
\newblock Local colourings and monochromatic partitions in complete
  bipartite graphs,
\newblock \emph{European J. Combin.} \textbf{60} (2017), 42--54,
\newblock \href{https://doi.org/10.1016/j.ejc.2016.09.003}{doi:10.1016/j.ejc.2016.09.003}.

\bibitem{Letzter2019}
S.~Letzter,
\newblock Monochromatic cycle partitions of $2$-coloured graphs with minimum
  degree $3n/4$,
\newblock \emph{Electron. J. Combin.} \textbf{26} (2019), Paper P1.19, 67~pp.,
\newblock \href{https://doi.org/10.37236/7239}{doi:10.37236/7239}.

\bibitem{LuczakRodlSzemeredi1998}
T.~\L{}uczak, V.~R\"odl and E.~Szemer\'edi,
\newblock Partitioning two-coloured complete graphs into two monochromatic
  cycles,
\newblock \emph{Combin. Probab. Comput.} \textbf{7} (1998), 423--436,
\newblock \href{https://doi.org/10.1017/S0963548398003599}{doi:10.1017/S0963548398003599}.

\bibitem{Pokrovskiy2014}
A.~Pokrovskiy,
\newblock Partitioning edge-coloured complete graphs into monochromatic
  cycles and paths,
\newblock \emph{J. Combin. Theory Ser. B} \textbf{106} (2014), 70--97,
\newblock \href{https://doi.org/10.1016/j.jctb.2014.01.003}{doi:10.1016/j.jctb.2014.01.003}.

\bibitem{Posa1963}
L.~P\'osa,
\newblock On the circuits of finite graphs,
\newblock \emph{Magyar Tud. Akad. Mat. Kutat\'o Int. K\"ozl.}
  \textbf{8} (1963), 355--361.

\bibitem{Sarkozy2011}
G.~N.~S\'ark\"ozy,
\newblock Monochromatic cycle partitions of edge-colored graphs,
\newblock \emph{J. Graph Theory} \textbf{66} (2011), 57--64,
\newblock \href{https://doi.org/10.1002/jgt.20492}{doi:10.1002/jgt.20492}.

\bibitem{Sarkozy2020}
G.~N.~S\'ark\"ozy,
\newblock Monochromatic partitions in local edge colorings,
\newblock \emph{Acta Math. Hungar.} \textbf{161} (2020), 412--426,
\newblock \href{https://doi.org/10.1007/s10474-020-01054-1}{doi:10.1007/s10474-020-01054-1}.

\end{thebibliography}

\end{document}
